\documentclass[../main]{subfiles}
\begin{document}
\ifSubfilesClassLoaded{\setcounter{page}{54}}{}
\section{Communism and Social Consciousness}
\label{sec:10_communism_consciousness}

From the standpoint of Marxism, in carrying out material production, society produces its own consciousness. Humanity has known various historically evolved forms of social consciousness. From mythological consciousness, in which the ideal is not yet separated, or, at various stages, not fully separated, from immediate material activity, philosophy and religion emerge with the development of the division of labour and the alienation of man from his own activity, followed by art and science. Alongside them exist morality and law, which emerged when society came to need external regulation of relations between people with opposing interests. Social consciousness is not the sum of the consciousnesses of separate individuals but the consciousness of society as a whole, separated from immediate human activity: its ideal dimension. Conversely, individual consciousness is that same social consciousness, acquired by a person according to his place and role in social life.

Social consciousness and its forms possess an objective existence in relation to the members of society. To live in human society, people are compelled to reckon with the demands and precepts of morality, law, religion, science, and even art and philosophy, just as they must reckon with surrounding nature and the level of technological development. What is reflected and expressed in social consciousness is not something arbitrary but the mode of interaction between man and nature and between man and man. Social consciousness and all its forms are not an immediate product of nature but a product of its assimilation by human society: a reflected mode of production.

The fact that social consciousness separates from immediate activity and that its forms acquire a relatively independent existence, as well as its compulsory, binding character in relation to the members of society, was recorded long before Marx. It underlies not only the religious worldview but all of idealist philosophy.

The becoming of the idea of communism required a critique both of this fact itself and of how it is interpreted in theoretical thinking. The critique of social consciousness, not merely of its own era but of any society in which the exploitation of man by man prevails, played a crucial role in the becoming of the ideological content of communism.

Marxism, as is well known, went through the school of criticism of Hegel's philosophy and of Feuerbach's critique of the Christian religion—a form of social consciousness that was dominant in Europe throughout the Middle Ages and continues to claim dominance in bourgeois society as well. Just how seriously Marx took Feuerbach's critique of religion can be seen from the fact that he speaks of nothing less than Feuerbach's great achievement. What, then, with respect to our question, does this achievement consist of, according to Marx? “(1) The proof that philosophy is nothing else but religion rendered into thought and expounded by thought, i.e., another form and manner of existence of the alienation of the human essence; hence equally to be condemned (this is a very important remark concerning philosophy as a form of social consciousness); (2) The establishment of true materialism and of real science, since Feuerbach also makes the social relation of ‘man to man’ the basic principle of the theory.”

These are the merits of Feuerbach's conception, but it is not entirely consistent, and Marx specifically pointed this out: “Feuerbach starts out from the fact of religious self-estrangement, of the duplication of the world into a religious, imaginary world and a real one. His work consists in resolving the religious world into its secular basis. He overlooks the fact that after completing this work, the chief thing still remains to be done. For the fact that the secular basis detaches itself from itself and establishes itself in the clouds as an independent realm can only be explained by the inner strife and intrinsic contradictoriness of this secular basis. The latter must itself, therefore, first be understood in its contradiction and then, by the removal of the contradiction, revolutionised in practice. Thus, for instance, once the earthly family is discovered to be the secret of the holy family, the former must then itself be criticised in theory and transformed in practice by revolutionary means.”

Man’s earthly life—material production—thus becomes for Marx the subject worth studying in order to grasp the essence of the earthly roots of consciousness. “Religion in itself is without content, it owes its being not to heaven but to the earth, and with the abolition of distorted reality, of which it is the theory, it will collapse of itself.” (Marx to Arnold Ruge in Dresden, 30 November 1842)

The same applies to the other forms of social consciousness: the already mentioned philosophy, as well as morality, law, and even science and art, must be understood as aspects of human practice—from practice itself, from its character and content, rather than from themselves. Marxism draws a clear distinction between the basis of social relations—the unity of the productive forces and the relations of production—and the superstructure derived from them, to which social consciousness belongs.

For example, in his preparatory manuscripts for Capital, Marx writes the following about law: “If these trivialities (the representation of law held by the economists contemporary to Marx—ed.) are reduced to their real content, they say more than their preachers know. Namely, that each form of production creates its own legal relations, form of government, etc. The crudity and lack of comprehension lies precisely in that organically coherent factors are brought into an accidental relation, into a merely reflective connection. All the bourgeois economists are aware of is that production can be carried on better under the modern police than, e.g., under club law. They forget only that club law is also law, and that the right of the strongest persists in another form even in their ‘rule-of-law state’.”

The “earthly roots” of social consciousness must not only be understood but also changed—the character of social production must be changed. And the fact that forms of social consciousness grow out of the mode of social production shows that matters with them are by no means as simple as they might seem at first glance. For example, law can be fully sublated, and the need for it as an external regulator overcome, only under definite social conditions—under communism developed on its own foundation. The same applies to overcoming religion. It cannot simply be abolished by declaring it false. It must be sublated, preserving all the achievements of human culture that developed in religious form, and done away with through a radical change in the character of the human activity carried out day after day.

Of course, in a small pamphlet on social consciousness under communism, we can speak only in general terms. But even so, one crucial circumstance must be noted concerning the class-divided society that precedes communism. Just as people occupy different positions in the production process according to their class affiliation, so too does social consciousness have a class character. This does not mean that there is no common field of consciousness, or no objective truth, goodness and beauty that do not depend solely and exclusively on the class position of their bearer. It means that this very field of consciousness is a field of class struggle. It means that the economically ruling class also rules in consciousness. That the consciousness of the ruling class is the ruling consciousness in class society. That is why Lenin, following Marx and Engels, speaks of the importance of struggle in the field of consciousness—of theoretical struggle as a special form of class struggle.

In his pamphlet “What Is to Be Done?”, Lenin devoted an entire section to the significance of the proletariat’s theoretical struggle and to Engels’s view on this question. The following remark by Lenin strikes us as highly relevant today: “Anyone who is at all acquainted with the actual state of our movement cannot but see that the wide spread of Marxism was accompanied by a certain lowering of the theoretical level. Quite a number of people with very little, and even a total lack of, theoretical training joined the movement because of its practical significance and its practical successes. We can judge, therefore, how tactless ‘Rabocheye Dyelo’ is when, with an air of triumph, it quotes Marx’s statement: ‘every step of real movement is more important than a dozen programmes.’ To repeat these words in a period of theoretical disorder is like shouting ‘many happy returns of the day!’ at the sight of a funeral procession. Moreover, these words of Marx are taken from his letter on the Gotha Programme, in which he sharply condemns the eclecticism displayed in the formulation of principles: if you must unite, Marx wrote to the party leaders, then enter into agreements to satisfy the practical aims of the movement, but do not allow any bargaining over principles, do not make theoretical ‘concessions.’ This was Marx’s idea, and yet there are people among us who, in his name, seek to belittle the significance of theory! Without revolutionary theory there can be no revolutionary movement. This idea cannot be insisted upon too strongly at a time when the fashionable preaching of opportunism goes hand in hand with an infatuation with the narrowest forms of practical activity.”

Turning to Engels’s legacy, Lenin cites a lengthy passage, which we consider it necessary to reproduce here as well: “The German workers have two important advantages over those of the rest of Europe. First, they belong to the most theoretical people of Europe, and they have retained that sense of theory which the so-called ‘educated’ classes of Germany have almost completely lost. Without German philosophy, which preceded it, particularly that of Hegel, German scientific socialism—the only scientific socialism that has ever existed—would never have come into being. Without a sense of theory among the workers, this scientific socialism would never have entered their flesh and blood as much as is the case today. What an immeasurable advantage this is may be seen, on the one hand, from the indifference towards all theory which is one of the main reasons why the English working-class movement crawls along so slowly in spite of the splendid organisation of the individual trades; and on the other hand, from the mischief and confusion wrought by Proudhonism, in its original form among the French and the Belgians, and in the form further caricatured by Bakunin among the Spaniards and the Italians.

The second advantage is that, chronologically speaking, the Germans were about the last to come into the workers’ movement. Just as German theoretical socialism will never forget that it rests on the shoulders of Saint-Simon, Fourier and Owen—three men who, for all their fantastic notions and all their utopianism, stand among the most eminent thinkers of all time, and whose genius anticipated countless truths whose correctness we are now proving scientifically—so the German practical workers’ movement must never forget that it has developed on the shoulders of the English and French movements, that it was able simply to turn to its own account their dearly bought experience and to avoid now the mistakes that, at the time, they for the most part could not avoid. Where would we be now without the example of the English trade unions and the French workers’ political struggles, without the colossal impetus given above all by the Paris Commune?

It must be said to the credit of the German workers that they have exploited the advantages of their situation with rare understanding. For the first time since a workers’ movement has existed, the struggle is being conducted pursuant to its three sides—the theoretical, the political and the practical-economic (resistance to the capitalists)—in harmony and in its interconnections, and in a systematic way.”

This passage shows Engels’s attitude towards theoretical struggle. But class struggle in the sphere of social consciousness is not exhausted by theoretical struggle alone, just as social consciousness is not confined to conceptual thinking. The struggle is also waged in the sphere of aesthetic activity, in the sphere of morality, and with respect to law and religion. This, too, must not be underestimated. Nevertheless, the theoretical work by which the proletariat acquires and develops its own class consciousness is the key to all other forms of consciousness, since it makes it possible to understand them.

It is also important to note that in class society, within social consciousness, the universally human is realised through the class-specific, whether or not people understand this. The consciousness of the revolutionary class, which at a given historical stage expresses the universally human interest, is both the driving force of history and the means by which this universally human content is realised. In this sense, not only morality, religion, philosophy and law, but also science and art are ideology as well, since both the aesthetic assimilation of the world in art and cognition in science are conditioned, and therefore limited (the freedom present in them is limited), by the historical movement of society and the balance of forces in the class struggle.

In this connection, I would like to turn to how Marx and Lenin resolved the problem of the subject of thinking. It should be recalled that Marx divided political economy into classical and vulgar, using a quite definite criterion to do so. He points out that classical political economy, unlike vulgar political economy, was a stage in the development of an understanding of its subject matter. He clearly defines the partisanship, that is, the class character, of this science. Both classical and vulgar political economy are bourgeois sciences, serving the bourgeoisie and expressing in theory the direction of its development. As soon as this development ceases to coincide with the development of society and ceases to embody that development—in other words, as soon as the bourgeoisie ceases to be a revolutionary class—political economy becomes vulgar. This means that it is no longer developing thought, but, on the contrary, becomes an obstacle to the development of thought.

Of course, the issue here is not the personal qualities of individual political economists. Marx convincingly shows that the subject of thinking is neither individual theorists nor society as a whole, but precisely a class. Moreover, a class is the subject of thinking only for as long as it is the subject of historical action.

We find the same idea in Lenin when he speaks of parties in philosophy, of which there can be only two: the party that thinks and the party that hinders thinking. And this is precisely because one party is the party of the class that is the subject of historical action, and therefore also the subject of thinking, while the other is not. This applies not only to philosophy but to ALL sciences without exception, including the natural sciences. Although it is less obvious there, it becomes so difficult, for example, to obtain funding for research that promises no returns even in the long run that such cases can quite reasonably be seen as episodes of a veiled class struggle. This remains relevant for as long as society develops through the struggle of classes. It follows that objective truth is always tendentious and has a class character. The same applies to the objectivity of goodness and beauty.

In September (August, Old Style) 1917, Lenin expressed a well-known and, in our view, very important idea in the spirit of what has been written above, one that shows how the proletariat relates to its class consciousness in revolutionary epochs, given that the aim of the proletariat, unlike that of other classes, is not at all to conquer and consolidate its own class domination, but, having conquered that domination, to abolish all class domination and every division of society into classes altogether: “Let us be unswerving,” says Lenin, “in submitting the slightest doubts to the judgement of the class-conscious workers, to the judgement of our Party. We trust it, we see in it the mind, honour and conscience of our epoch, and in the international alliance of revolutionary internationalists we see the sole guarantee of the liberation movement of the working class.” “The domination of a class is impossible without its domination in the sphere of consciousness,” Lenin emphasises. In order to win, the revolutionary class must, in the sphere of consciousness, measure itself not by someone else’s yardstick but by its own, one determined by its historical mission. Having wrested domination in this sphere from the bourgeoisie and armed itself with the most advanced achievements of its own class consciousness, the proletariat becomes a genuine material force of historical transformation: “The weapon of criticism cannot, of course, replace criticism by weapons, material force must be overthrown by material force; but theory also becomes a material force as soon as it has gripped the masses,” Marx wrote in “Contribution to the Critique of Hegel’s Philosophy of Law”.

So far, however, we have been speaking not of social consciousness under communism, but of the state of affairs from which we must start in the transition to it. What, then, does communism as a science insist upon with regard to social consciousness?

Communism insists on overcoming the social division of labour and the alienation of human beings from other people and from the process and the result of their own activity. And this entails corresponding changes in the sphere of social consciousness as well. The alienation of the forms of social consciousness, both from one another and directly from human activity, will also be overcome.

The ideal active capacities of human beings, which are “locked up” within the forms of social consciousness as spheres of aesthetic, moral and intellectual activity separated from material production, will become the possession of each and every one.

The question of social consciousness under communism is a question of practice. It is a question of planning, and of carrying out what has been planned in material activity, in such a way that corresponding changes also take place in the spiritual sphere. By changing their activity, people also change their consciousness as the ideal side of that activity. But under communism, the very change in practice is conscious in character. That is why here it is no longer simply that social being determines social consciousness; rather, social consciousness, as the consciousness of a free association of people acting in concert, determines the being that determines this consciousness.

The socialisation of all spiritual culture, and moreover in accordance with the logic of that culture, becomes a most acute necessity for the transition to communism. That is why, in the theory of communism too, the problems of social consciousness make sense only if they are understood as questions of the becoming of communism.

In communist activity, unlike in alienated labour, the ideal as that which corresponds to an ideal and the ideal as that which pertains to the “world of ideas” coincide empirically in the last instance (precisely in the last instance). The ideal thus appears at once as the reflection (representation, expression) of the material in the material—an active reflection of objective reality—as the ideal goal of activity, to which that activity must conform, and as an aesthetic ideal. Herein lies the most important moment in overcoming the division of labour and its alienation in every sense of the word, and the intrinsic value of activity is affirmed regardless of what exactly is being done.

It should be noted here that it would be wrong to claim that the conscious and material moments of human activity are on an equal footing in labour as such, regardless of its historical specificity, let alone at the level of the activity of society as a whole. Of course, the ideal moment is present in any activity as the moment of goal-setting, unless goal-setting lies outside that activity (for a worker who merely sells his labour power, his own goal does not coincide with the goal of production at all), but that is not the point. In any case, the mode of production of material life determines the mode of production of the ideal. And this has to be reckoned with.

This matters in practical terms because, although an idea does become a material force once it has gripped the masses, the idea itself is the result of the conscious grasp of material forces in the course of their action (the result and an aspect of the interaction of the material with the material). In other words, to change the mode of production of ideas, one must change the mode of production of material life. Only this way, and not the other way round. The other way round, nothing will come of it.

For example, however much we propagandise dialectics today (that is, act ideally), it will not thereby become everyday thinking, and the production of dialectical thinking will for now remain a cottage industry, so to speak. This cottage industry is, of course, very important, and much of the future depends on it. But we must understand that dialectically thinking people will not be a mass product under the present mode of production of material life and the mode of production of consciousness that corresponds to it.

For the problem of communism, it is not enough simply to establish the direction in which man’s material life passes into consciousness (what is primary and what is secondary); one must also put forward ideas for practically accomplishing this transition. And here particular ideas on this score are very important. Ideas for shaping various empirical forms of communism are important, since the material never passes into consciousness all at once and in its entirety, but always does so in certain particular forms. Ideas such as Makarenko’s, for example, are important: love is an ordinary matter that must be properly organised.

However enormous the role played by the ideal moment of activity, what is always at issue is the life and activity of material man in the material world. The ideal is always an attribute of the material activity of social man and corresponds to its movement. The “realm of freedom”, total communism, is characterised, among other things, by the conscious production of social consciousness on a scale that is likewise social.

Isolated breakthroughs in the sphere of social consciousness that go beyond the bounds of capitalism are already taking place. They were discussed (among other things) in the second chapter of this work. This is the production of the ideal in the creation of modern technologies. But that is not the whole story: not all of the ideal belongs to social consciousness. This is no longer simply ideal activity in its alienation, which is superstructural in character, nor is it the ideal confined within the limits of professional cretinism. It is the ideal moment of the production of human life in general. And precisely how this life is produced determines (as with databases and free software, one and the same process is at work here) how a new, immediate, unalienated social consciousness is created, even if only in local areas for now. This is a consciousness produced not only, and not so much, by the mere fact of producing a program as by the very character of its production and functioning.

Social consciousness can determine social consciousness only by determining social being. Moreover, the latter must not be understood as something separate from the whole human being, as something ideal in general, but only as a moment of the activity of the material human being in the equally material world. The point lies precisely in the conscious production of the whole human being: one who acts, feels and thinks universally. Such activity, regardless of what exactly is being done, will appear as the realisation of communism in a directly collective form and will be activity aimed at producing freedom, goal-setting, sensuousness, morality and thinking.

By and large, this presupposes that any activity under communism must, among other things, be the production of ideas and of the ideal of future life, the production of the human being.
\end{document}
